\documentclass[14pt,a4paper]{extarticle}
% =============================================================
% preamble.tex — оформление отчётов по СТП БГУИР 01-2024
% =============================================================

% Шрифты: под pdflatex — Times через tempora + T2A,
% под xelatex/tectonic — системный Times New Roman (запасной вариант STIX Two Text)
\usepackage{iftex}
\ifPDFTeX
  \usepackage[utf8]{inputenc}
  \usepackage[T2A]{fontenc}
  \usepackage{tempora}
\else
  \usepackage{fontspec}
  \IfFontExistsTF{Times New Roman}
    {\setmainfont{Times New Roman}}
    {\setmainfont{STIX Two Text}}
  \IfFontExistsTF{DejaVu Sans Mono}
    {\setmonofont{DejaVu Sans Mono}[Scale=0.82]}
    {\setmonofont{DejaVuSansMono}[Extension=.ttf, BoldFont=*-Bold,
                                  ItalicFont=*-Oblique, Scale=0.82]}
\fi
\usepackage[russian]{babel}

% Поля: слева 30 мм, справа 15 мм, сверху и снизу 20 мм
\usepackage{geometry}
\geometry{left=30mm, right=15mm, top=20mm, bottom=20mm}

\usepackage{setspace}
\singlespacing

\pretolerance=1000
\tolerance=2000
\emergencystretch=3em

\usepackage{indentfirst}
\setlength{\parindent}{1.25cm}

% Номер страницы — снизу справа
\usepackage{fancyhdr}
\pagestyle{fancy}
\fancyhf{}
\fancyfoot[R]{\thepage}
\renewcommand{\headrulewidth}{0pt}
\fancypagestyle{plain}{%
  \fancyhf{}%
  \fancyfoot[R]{\thepage}%
  \renewcommand{\headrulewidth}{0pt}%
}

\usepackage{amsmath,amssymb}
\usepackage{graphicx}
\usepackage{float}
\usepackage{booktabs}
\usepackage{tabularx}
\usepackage{multirow}
\usepackage{longtable}
\usepackage{array}

\usepackage{caption}
\captionsetup[figure]{labelsep=endash, justification=centering,
                      font={normalsize}, name=Рисунок}
\captionsetup[table]{labelsep=endash, justification=raggedright,
                     singlelinecheck=false, font={normalsize}, name=Таблица}

\usepackage{titlesec}
\titleformat{\section}
  {\newpage\normalfont\fontsize{14}{18}\selectfont\bfseries}
  {\thesection}{0.5em}{\MakeUppercase}
\titlespacing*{\section}{1.25cm}{0pt}{14pt}
\titleformat{\subsection}
  {\normalfont\fontsize{14}{18}\selectfont\bfseries}
  {\thesubsection}{0.5em}{}
\titlespacing*{\subsection}{1.25cm}{14pt}{14pt}
\titleformat{\subsubsection}
  {\normalfont\fontsize{14}{18}\selectfont\bfseries}
  {\thesubsubsection}{0.5em}{}
\titlespacing*{\subsubsection}{1.25cm}{14pt}{14pt}

\newcommand{\sectioncenter}[1]{%
  \newpage
  \phantomsection
  \addcontentsline{toc}{section}{#1}
  \begin{center}\textbf{\MakeUppercase{#1}}\end{center}
  \vspace{14pt}
}

\usepackage{tocloft}
\renewcommand{\cftsecfont}{\normalfont}
\renewcommand{\cftsecpagefont}{\normalfont}
\renewcommand{\cftsubsecfont}{\normalfont}
\renewcommand{\cftsubsecpagefont}{\normalfont}
\renewcommand{\cftsecleader}{\cftdotfill{\cftdotsep}}
\renewcommand{\cftdotsep}{1}
% babel переопределяет заголовки после преамбулы, поэтому правим через \addto
\addto\captionsrussian{\renewcommand{\contentsname}{\hfill\textbf{СОДЕРЖАНИЕ}\hfill}}
\renewcommand{\cftaftertoctitle}{\hfill}

\usepackage{enumitem}
\setlist[itemize]{noitemsep, topsep=0pt, leftmargin=1.25cm, label=--}
\setlist[enumerate]{noitemsep, topsep=0pt, leftmargin=1.25cm}

% Листинги кода и вывода программы
\usepackage{listings}
\usepackage{fancyvrb}
\lstset{
  basicstyle=\ttfamily\footnotesize,
  frame=single,
  breaklines=true,
  showstringspaces=false,
  inputencoding=utf8,
  extendedchars=true,
  columns=fullflexible,
  keepspaces=true,
  captionpos=t
}
\ifPDFTeX
\lstset{
  literate={а}{{\cyra}}1 {б}{{\cyrb}}1 {в}{{\cyrv}}1 {г}{{\cyrg}}1
           {д}{{\cyrd}}1 {е}{{\cyre}}1 {ж}{{\cyrzh}}1 {з}{{\cyrz}}1
           {и}{{\cyri}}1 {й}{{\cyrishrt}}1 {к}{{\cyrk}}1 {л}{{\cyrl}}1
           {м}{{\cyrm}}1 {н}{{\cyrn}}1 {о}{{\cyro}}1 {п}{{\cyrp}}1
           {р}{{\cyrr}}1 {с}{{\cyrs}}1 {т}{{\cyrt}}1 {у}{{\cyru}}1
           {ф}{{\cyrf}}1 {х}{{\cyrh}}1 {ц}{{\cyrc}}1 {ч}{{\cyrch}}1
           {ш}{{\cyrsh}}1 {щ}{{\cyrshch}}1 {ъ}{{\cyrhrdsn}}1
           {ы}{{\cyrery}}1 {ь}{{\cyrsftsn}}1 {э}{{\cyrerev}}1
           {ю}{{\cyryu}}1 {я}{{\cyrya}}1 {ё}{{\cyryo}}1
           {А}{{\CYRA}}1 {Б}{{\CYRB}}1 {В}{{\CYRV}}1 {Г}{{\CYRG}}1
           {Д}{{\CYRD}}1 {Е}{{\CYRE}}1 {Ж}{{\CYRZH}}1 {З}{{\CYRZ}}1
           {И}{{\CYRI}}1 {Й}{{\CYRISHRT}}1 {К}{{\CYRK}}1 {Л}{{\CYRL}}1
           {М}{{\CYRM}}1 {Н}{{\CYRN}}1 {О}{{\CYRO}}1 {П}{{\CYRP}}1
           {Р}{{\CYRR}}1 {С}{{\CYRS}}1 {Т}{{\CYRT}}1 {У}{{\CYRU}}1
           {Ф}{{\CYRF}}1 {Х}{{\CYRH}}1 {Ц}{{\CYRC}}1 {Ч}{{\CYRCH}}1
           {Ш}{{\CYRSH}}1 {Щ}{{\CYRSHCH}}1 {Ъ}{{\CYRHRDSN}}1
           {Ы}{{\CYRERY}}1 {Ь}{{\CYRSFTSN}}1 {Э}{{\CYREREV}}1
           {Ю}{{\CYRYU}}1 {Я}{{\CYRYA}}1 {Ё}{{\CYRYO}}1
           {—}{{---}}1 {–}{{--}}1 {…}{{\ldots}}1
}
\fi

% Подписи листингов — в том же стиле, что таблицы: «Листинг X.Y — Название»
\captionsetup[lstlisting]{labelsep=endash, justification=raggedright,
                          singlelinecheck=false, font={normalsize}}
\renewcommand{\lstlistingname}{Листинг}
\AtBeginDocument{\numberwithin{lstlisting}{section}}

\lstdefinestyle{python}{language=Python, keywordstyle=\bfseries}
\lstdefinestyle{shell}{language=bash, keywordstyle=\mdseries}
\lstdefinestyle{plain}{language={}, basicstyle=\ttfamily\scriptsize}

% Абзац после листинга должен сохранять отступ 1,25 см (СТП 2.1.1)
\makeatletter
\AtBeginDocument{\let\@endpetrue\@endpefalse}
\makeatother

% Подпись к блоку verbatim: \captionof вне плавающего окружения
% глобально обнуляет \parindent, поэтому отступ восстанавливается вручную
\newcommand{\listingcaption}[2]{%
  \captionof{lstlisting}{#1}\label{#2}%
  \setlength{\parindent}{1.25cm}%
}

\usepackage{hyperref}
\hypersetup{colorlinks=true, linkcolor=black, urlcolor=blue, citecolor=black}

\numberwithin{figure}{section}
\numberwithin{table}{section}
\numberwithin{equation}{section}

% Титульный лист отчёта по лабораторной работе
\newcommand{\makelabtitle}[2]{%
  \begin{titlepage}
  \thispagestyle{empty}
  \begin{center}

  МИНИСТЕРСТВО ОБРАЗОВАНИЯ РЕСПУБЛИКИ БЕЛАРУСЬ

  \vspace{0.5em}
  Учреждение образования

  \vspace{0.5em}
  \textbf{<<Белорусский государственный университет информатики и радиоэлектроники>>}

  \vspace{2em}
  Факультет компьютерных систем и сетей

  \vspace{0.5em}
  Кафедра программного обеспечения информационных технологий
  \end{center}

  \vspace{3em}
  \begin{center}
  \textbf{ОТЧЕТ}

  \vspace{0.5em}
  по лабораторной работе №#1

  \vspace{0.5em}
  по дисциплине

  \vspace{0.5em}
  <<Информационные и технические средства защиты авторских прав>>

  \vspace{0.5em}
  <<#2>>
  \end{center}

  \vspace{5em}
  \begin{flushright}
  \begin{minipage}{0.5\textwidth}
  \textbf{Выполнил:}\\
  Лебедевич Артем Владимирович\\
  магистрант группы 556301

  \vspace{0.6em}
  \textbf{Преподаватель:}\\
  Ярмолик Вячеслав Николаевич\\
  доктор технических наук\\
  профессор кафедры ПОИТ
  \end{minipage}
  \end{flushright}

  \vfill
  \begin{center}Минск, 2026~г.\end{center}
  \end{titlepage}
}

% Титульный лист реферата
\newcommand{\makerefattitle}[1]{%
  \begin{titlepage}
  \thispagestyle{empty}
  \begin{center}

  МИНИСТЕРСТВО ОБРАЗОВАНИЯ РЕСПУБЛИКИ БЕЛАРУСЬ

  \vspace{0.5em}
  Учреждение образования

  \vspace{0.5em}
  \textbf{<<Белорусский государственный университет информатики и радиоэлектроники>>}

  \vspace{2em}
  Факультет компьютерных систем и сетей

  \vspace{0.5em}
  Кафедра программного обеспечения информационных технологий
  \end{center}

  \vspace{3em}
  \begin{center}
  \textbf{РЕФЕРАТ}

  \vspace{0.5em}
  по дисциплине

  \vspace{0.5em}
  <<Информационные и технические средства защиты авторских прав>>

  \vspace{1.5em}
  на тему

  \vspace{0.5em}
  \textbf{<<#1>>}
  \end{center}

  \vspace{5em}
  \begin{flushright}
  \begin{minipage}{0.5\textwidth}
  \textbf{Выполнил:}\\
  Лебедевич Артем Владимирович\\
  магистрант группы 556301

  \vspace{0.6em}
  \textbf{Преподаватель:}\\
  Ярмолик Вячеслав Николаевич\\
  доктор технических наук\\
  профессор кафедры ПОИТ
  \end{minipage}
  \end{flushright}

  \vfill
  \begin{center}Минск, 2026~г.\end{center}
  \end{titlepage}
}

\usepackage{tikz}
\usetikzlibrary{arrows.meta, positioning}

\begin{document}

\makerefattitle{Доказательство с нулевым знанием}

\tableofcontents

\sectioncenter{Введение}

Классическая криптография отвечает на вопрос, как передать сообщение так, чтобы его не
прочитал посторонний. Однако во множестве практических задач требуется решить другую,
на первый взгляд парадоксальную задачу: убедить собеседника в истинности утверждения,
не сообщив ему ничего сверх самого факта истинности. Владелец банковской карты должен
доказать терминалу, что знает секретный ключ, но не передавать этот ключ; пользователь
сервиса~--- что ему исполнилось 18 лет, не раскрывая даты рождения; оператор блокчейна~---
что транзакция корректна, не публикуя ни отправителя, ни суммы.

Формальный аппарат для таких задач~--- \emph{доказательства с нулевым знанием}
(zero-knowledge proofs, ZKP). Понятие введено в 1985~году в работе Голдвассер, Микали и
Ракоффа <<The Knowledge Complexity of Interactive Proof-Systems>>, за которую авторы
впоследствии получили премию Гёделя, а Голдвассер и Микали~--- премию Тьюринга.
Первые практические протоколы (Фиата~--- Шамира, Гиллу~--- Кискатра, Шнорра)
появились уже во второй половине 1980-х годов и предназначались для идентификации
владельцев смарт-карт при жёстких ограничениях на память и вычислительную мощность.

За последнее десятилетие тема пережила второе рождение: конструкции zk-SNARK и
zk-STARK позволяют доказывать корректность выполнения произвольной программы
доказательством в несколько сотен байт, проверяемым за миллисекунды. На этой основе
работают приватные платёжные системы, механизмы масштабирования блокчейнов
(zk-rollups), системы анонимной аутентификации и подтверждения происхождения контента.

Цель реферата~--- рассмотреть доказательства с нулевым знанием от формального
определения и классических протоколов идентификации, изучаемых в курсе, до современных
конструкций и областей их практического применения, а также оценить ограничения и
риски внедрения этой технологии.

\section{Понятие доказательства с нулевым знанием}

\subsection{Интерактивная система доказательства}

Рассматривается взаимодействие двух сторон: \emph{доказывающего}~$P$ (prover), который
обладает секретом и стремится убедить собеседника, и \emph{проверяющего}~$V$ (verifier),
который должен принять или отвергнуть доказательство. Обе стороны~--- вероятностные
алгоритмы, обменивающиеся сообщениями; $P$ в теоретической модели не ограничен по
вычислительным ресурсам, $V$ работает за полиномиальное время.

Пусть $L$~--- некоторый язык (множество истинных утверждений), а $x$~--- утверждение,
истинность которого доказывается, например: <<число~$v$ является квадратичным вычетом по
модулю~$n$>> или <<мне известен ключ, соответствующий данному открытому ключу>>. Пара
$(P, V)$ образует интерактивную систему доказательства для~$L$, если выполняются два
первых свойства из перечисленных ниже, и систему \emph{с нулевым знанием}~--- если
выполняется также третье.

\subsection{Три свойства}

\begin{enumerate}
\item \textbf{Полнота} (completeness). Если утверждение истинно и обе стороны действуют
по протоколу, то проверяющий принимает доказательство с вероятностью, близкой к единице.
\item \textbf{Корректность} (soundness). Если утверждение ложно, то никакой,
в том числе нечестный, доказывающий не сможет убедить проверяющего с вероятностью выше
заданной малой величины~$\varepsilon$. Величина $\varepsilon$ снижается повторением
раундов протокола: для базового раунда с вероятностью обмана~$1/2$ после $t$ независимых
раундов она составляет $2^{-t}$.
\item \textbf{Нулевое разглашение} (zero-knowledge). Проверяющий не узнаёт из протокола
ничего, кроме самого факта истинности утверждения.
\end{enumerate}

Третье свойство формализуется через \emph{симулятор}: система обладает нулевым
разглашением, если существует вероятностный полиномиальный алгоритм~$S$, который без
доступа к секрету порождает протокольные записи (транскрипты), неотличимые от
записей реального взаимодействия. Раз такие записи можно изготовить самостоятельно,
они не несут проверяющему никакой дополнительной информации. В зависимости от того,
насколько строго понимается неотличимость, различают совершенное, статистическое и
вычислительное нулевое разглашение.

Отдельно различают \emph{доказательства} (proofs), корректность которых сохраняется
против неограниченного по ресурсам противника, и \emph{аргументы} (arguments),
корректные лишь против полиномиально ограниченного противника. Практически все
современные компактные конструкции, включая zk-SNARK, являются аргументами.

\subsection{Пещера Али-Бабы}

Классическая иллюстрация, предложенная Кискатром и соавторами, описывает кольцевую
пещеру с развилкой: два коридора~$A$ и~$B$ соединены дверью, открывающейся только по
секретному слову. Пегги (доказывающий) утверждает, что знает слово, Виктор
(проверяющий) хочет в этом убедиться, не узнав слова.

Протокол одного раунда: Пегги входит в пещеру и выбирает коридор, Виктор подходит к
развилке и называет коридор, из которого Пегги должна выйти. Если слово ей известно,
она выходит из любого названного коридора; если нет~--- угадывает нужный коридор с
вероятностью~$1/2$. После $t$ раундов вероятность обмана падает до $2^{-t}$, при этом
Виктор не получает ни бита информации о самом слове.

Существенно, что нулевое разглашение здесь опирается на \emph{случайность вызова}:
запись протокола, сделанная посторонним наблюдателем, не убеждает никого, поскольку
её легко подделать, договорившись о последовательности коридоров заранее. Это же
свойство отличает интерактивное доказательство от цифровой подписи, которая, наоборот,
должна убеждать любую третью сторону.

\subsection{Структура $\Sigma$-протокола}

Все рассматриваемые далее классические протоколы построены по единой трёхшаговой
схеме, которую называют $\Sigma$-протоколом (рисунок~\ref{fig:sigma}):

\begin{enumerate}
\item \textbf{обязательство} (commitment): $P$ выбирает случайное значение и посылает
производную от него величину~$x$, не раскрывая самой случайности;
\item \textbf{вызов} (challenge): $V$ посылает случайный вызов~$e$, выбранный из
заранее оговорённого множества;
\item \textbf{ответ} (response): $P$ вычисляет~$y$, связывающий случайность, секрет и
вызов, а $V$ проверяет соотношение между $x$, $e$ и~$y$.
\end{enumerate}

\begin{figure}[H]
\centering
\begin{tikzpicture}[font=\small, >={Stealth[length=3mm]}]
  \node (p) at (0,0) {\textbf{Доказывающий $P$}};
  \node (v) at (9,0) {\textbf{Проверяющий $V$}};
  \node[font=\footnotesize, below=2pt of p] {знает секрет $s$};
  \node[font=\footnotesize, below=2pt of v] {знает открытый ключ $v$};

  \draw[thick] (0,-1.2) -- (0,-4.3);
  \draw[thick] (9,-1.2) -- (9,-4.3);

  \draw[->, thick] (0,-1.8) -- node[above]{обязательство $x$} (9,-1.8);
  \draw[<-, thick] (0,-2.8) -- node[above]{вызов $e$ (случайный)} (9,-2.8);
  \draw[->, thick] (0,-3.8) -- node[above]{ответ $y = f(r, s, e)$} (9,-3.8);

  \node[font=\footnotesize, anchor=north east] at (9.6,-4.4)
    {проверка $g(x, e, y) = \text{истина}$};
\end{tikzpicture}
\caption{Общая схема $\Sigma$-протокола}
\label{fig:sigma}
\end{figure}

Корректность таких протоколов обычно доказывается через свойство \emph{специальной
корректности}: из двух принимаемых ответов на разные вызовы при одном и том же
обязательстве можно вычислить секрет. Отсюда следует, что доказывающий, способный
отвечать более чем на один вызов, секретом действительно владеет.

\section{Классические протоколы идентификации}

\subsection{Протокол Фиата~--- Шамира}

Протокол предложен А.~Фиатом и А.~Шамиром в 1986~году и основан на трудности извлечения
квадратного корня по составному модулю, что эквивалентно задаче факторизации.

\textbf{Подготовка.} Доверенный центр выбирает модуль $n = p \cdot q$, где $p$ и $q$~---
большие простые числа, известные только центру. Пользователь выбирает секрет~$s$,
взаимно простой с~$n$, и публикует
\begin{equation}
v = s^2 \bmod n.
\end{equation}

\textbf{Раунд протокола.}
\begin{enumerate}
\item $P$ выбирает случайное $r \in [1, n)$ и посылает $x = r^2 \bmod n$;
\item $V$ посылает случайный бит $e \in \{0, 1\}$;
\item $P$ отвечает $y = r \cdot s^{\,e} \bmod n$;
\item $V$ проверяет равенство $y^2 \equiv x \cdot v^{\,e} \pmod n$.
\end{enumerate}

При $e = 0$ проверяется $y^2 \equiv x$, то есть $P$ демонстрирует знание корня из
обязательства; при $e = 1$~--- $y^2 \equiv x \cdot v$, что требует знания~$s$. Обманщик
может подготовиться к одному из двух вызовов, но не к обоим сразу: угадав заранее $e=1$,
он выбирает $y$ произвольно и посылает $x = y^2 v^{-1}$. Поэтому вероятность обмана в
одном раунде равна~$1/2$, а после $t$ раундов~--- $2^{-t}$.

Нулевое разглашение видно из работы симулятора: не зная~$s$, он угадывает вызов~$e$,
подбирает $y$ случайно и вычисляет $x = y^2 v^{-e} \bmod n$; при неверном угадывании
попытка отбрасывается. Полученные транскрипты распределены так же, как настоящие.

В обобщении Фейге~--- Фиата~--- Шамира пользователь владеет $k$ секретами
$s_1, \dots, s_k$ с открытыми значениями $v_i = s_i^2 \bmod n$, вызов представляет
собой вектор битов $(e_1, \dots, e_k)$, а ответ имеет вид
$y = r \prod_{i=1}^{k} s_i^{\,e_i} \bmod n$. Вероятность обмана за $t$ раундов
снижается до $2^{-kt}$, что позволяет сократить число раундов за счёт увеличения
объёма хранимого секрета.

\subsection{Протокол Гиллу~--- Кискатра}

Протокол, предложенный Л.~Гиллу и Ж.-Ж.~Кискатром в 1988~году, разрабатывался
непосредственно для смарт-карт и снижает число раундов до одного за счёт расширения
множества вызовов. Его стойкость опирается на трудность задачи RSA.

\textbf{Подготовка.} Центр публикует модуль~$n$ и открытую экспоненту~$v$. Идентификатор
пользователя отображается в число~$J$, а центр выдаёт пользователю секрет~$B$, такой что
\begin{equation}
J \cdot B^{\,v} \equiv 1 \pmod n .
\end{equation}

\textbf{Протокол.}
\begin{enumerate}
\item $P$ выбирает случайное~$r$ и посылает $T = r^{\,v} \bmod n$;
\item $V$ посылает вызов $d \in [0, v)$;
\item $P$ отвечает $D = r \cdot B^{\,d} \bmod n$;
\item $V$ проверяет $T \equiv D^{\,v} \cdot J^{\,d} \pmod n$.
\end{enumerate}

Вероятность обмана в одном раунде составляет $1/v$, поэтому при достаточно большой
экспоненте (например, $v \approx 2^{20}$) хватает единственного раунда. Это принципиально
для смарт-карт: каждый раунд~--- это обмен сообщениями с терминалом и несколько
возведений в степень по большому модулю. Протоколы данного семейства вошли в
международный стандарт ISO/IEC~9798-5, описывающий механизмы аутентификации
с нулевым разглашением.

\subsection{Протокол Шнорра}

К.~Шнорр предложил в 1989~году протокол, основанный на трудности дискретного
логарифмирования в подгруппе простого порядка, наиболее эффективный из классических.

\textbf{Подготовка.} Выбираются простые $p$ и $q$ такие, что $q \mid p - 1$, и элемент~$g$
порядка~$q$ по модулю~$p$. Секрет~--- случайное $a \in [1, q)$, открытый ключ~---
\begin{equation}
v = g^{-a} \bmod p .
\end{equation}

\textbf{Протокол.}
\begin{enumerate}
\item $P$ выбирает случайное $r \in [1, q)$ и посылает $x = g^{\,r} \bmod p$;
\item $V$ посылает вызов $e \in [0, 2^{t})$;
\item $P$ отвечает $y = r + a \cdot e \bmod q$;
\item $V$ проверяет $x \equiv g^{\,y} \cdot v^{\,e} \pmod p$.
\end{enumerate}

Вероятность обмана равна $2^{-t}$ за один проход. Протокол требует от доказывающего
одного возведения в степень, причём обязательство $x$ не зависит от вызова и может быть
вычислено заранее, в том числе в свободное время смарт-карты. Именно эта схема легла в
основу современных цифровых подписей~--- подписи Шнорра, DSA и EdDSA.

\subsection{Сравнение классических протоколов}

\begin{table}[H]
\caption{Классические протоколы идентификации с нулевым знанием}
\label{tab:classic}
\centering
\small
\setlength{\tabcolsep}{5pt}
\begin{tabularx}{\textwidth}{>{\raggedright\arraybackslash}p{3.8cm}X>{\raggedright\arraybackslash}p{3.5cm}c}
\toprule
Протокол & Основание стойкости & Вероятность обмана & Год \\
\midrule
Фиата~--- Шамира & извлечение квадратного корня по составному модулю & $2^{-t}$ за $t$ раундов & 1986 \\
Фейге~--- Фиата~--- Шамира & то же, $k$ секретов & $2^{-kt}$ за $t$ раундов & 1988 \\
Гиллу~--- Кискатра & задача RSA & $v^{-t}$, достаточно $t = 1$ & 1988 \\
Шнорра & дискретный логарифм в подгруппе порядка $q$ & $2^{-t}$ за один проход & 1989 \\
\bottomrule
\end{tabularx}
\end{table}

Общая тенденция очевидна: развитие шло по пути сокращения числа раундов и объёма
пересылаемых данных при сохранении той же вероятности обмана. Практическим следствием
стало применение протоколов в банковских картах, системах контроля доступа и
телевизионных смарт-картах условного доступа.

\section{От интерактивного протокола к цифровой подписи}

\subsection{Преобразование Фиата~--- Шамира}

Интерактивность~--- главное практическое неудобство классических протоколов: стороны
должны быть на связи одновременно, а полученный транскрипт не убеждает третью сторону.
Преобразование Фиата~--- Шамира устраняет оба недостатка: случайный вызов проверяющего
заменяется значением криптографической хеш-функции от обязательства (и, при
необходимости, от подписываемого сообщения):
\begin{equation}
e = H(x \parallel m).
\end{equation}

Доказывающий не может подобрать выгодный ему вызов, поскольку не контролирует выход
хеш-функции, а любой желающий может пересчитать~$e$ и проверить соотношение. Протокол
из трёхшагового превращается в одно сообщение~$(x, y)$, то есть в
\emph{неинтерактивное} доказательство, а если в хеш включено сообщение~$m$~--- в
цифровую подпись.

Именно так из протокола Шнорра получается подпись Шнорра, а из неё~--- схема EdDSA
(Ed25519), применяемая сегодня в SSH, TLS, подписи пакетов и блокчейн-системах.
Аналогично из протокола Гиллу~--- Кискатра получается подпись GQ. Таким образом,
раздел курса о доказательствах с нулевым знанием напрямую связан с разделом о цифровой
подписи: это две стороны одной конструкции.

Строгое обоснование безопасности такого преобразования выполняется в \emph{модели
случайного оракула}, где хеш-функция считается идеально случайной. Модель является
идеализацией, но на практике преобразование Фиата~--- Шамира считается надёжным при
корректной реализации.

\subsection{Типичная ошибка реализации}

Практическая опасность преобразования состоит в том, что в хеш необходимо включать
\emph{все} публичные параметры доказательства, а не только обязательство. Если часть
данных в хеш не попадает, нечестный доказывающий получает свободу выбора и может
изготовить доказательство ложного утверждения~--- это называют <<слабым>>
преобразованием Фиата~--- Шамира.

В 2022~году исследователи Trail of Bits описали серию уязвимостей такого рода под общим
названием Frozen Heart, обнаружив их сразу в нескольких широко используемых
библиотеках доказательств. Особенность подобных ошибок в том, что они не проявляются
при обычном тестировании: честно построенные доказательства проверяются успешно, и лишь
целенаправленная атака выявляет возможность подделки.

\section{Современные конструкции}

Классические протоколы доказывают знание конкретного числа. Современные конструкции
решают несравнимо более общую задачу: доказать, что доказывающему известны такие
входные данные~$w$, при которых заданная программа~$C$ возвращает истину,
$C(x, w) = 1$, где $x$~--- публичные параметры. Программа предварительно
переводится в систему полиномиальных или арифметических ограничений (арифметизация),
после чего доказательство строится уже для неё.

\subsection{zk-SNARK}

Аббревиатура SNARK означает <<succinct non-interactive argument of knowledge>>~---
краткий неинтерактивный аргумент знания. Ключевое свойство~--- \emph{краткость}: размер
доказательства и время проверки практически не зависят от объёма вычисления, которое
доказывается. В схеме Groth16 (2016) доказательство состоит из трёх элементов
эллиптической группы, то есть занимает порядка двух сотен байт, а проверка сводится к
нескольким операциям спаривания и выполняется за единицы миллисекунд.

Плата за краткость~--- \emph{доверенная установка} (trusted setup): перед началом работы
необходимо сгенерировать открытые параметры, при генерации которых возникают секретные
величины. Если они не уничтожены, их владелец способен создавать доказательства ложных
утверждений (<<токсичные отходы>> установки). На практике это решается церемониями с
множеством независимых участников, где достаточно честности хотя бы одного, а также
схемами с универсальной обновляемой установкой (PLONK, 2019), где одна церемония
обслуживает любые программы заданного размера.

\subsection{zk-STARK}

Конструкция STARK (Бен-Сассон и соавторы, 2018) отказывается от доверенной установки:
она \emph{прозрачна} (transparent), то есть использует только публичную случайность.
Стойкость опирается лишь на криптографические хеш-функции, без предположений о
сложности задач теории чисел, поэтому схема считается устойчивой к атакам квантового
компьютера~--- в отличие от конструкций на эллиптических кривых, уязвимых к алгоритму
Шора.

Обратная сторона~--- размер доказательства: десятки и сотни килобайт против сотен байт
у Groth16. Поэтому выбор между SNARK и STARK~--- это выбор между стоимостью хранения и
публикации доказательства и требованиями к доверию и постквантовой стойкости.

\subsection{Bulletproofs}

Bulletproofs (Бюнц и соавторы, 2018) занимают промежуточное положение: доверенная
установка не требуется, размер доказательства растёт логарифмически от размера
утверждения. Конструкция особенно эффективна для \emph{доказательств диапазона}
(range proofs)~--- утверждений вида <<скрытое значение лежит в интервале
$[0, 2^{64})$>>, необходимых в конфиденциальных транзакциях, чтобы исключить
отрицательные суммы. Слабое место~--- время проверки, растущее линейно, что ограничивает
применение в системах с массовой верификацией.

\subsection{Сравнение современных конструкций}

\begin{table}[H]
\caption{Сравнение современных систем доказательств}
\label{tab:modern}
\centering
\small
\setlength{\tabcolsep}{4pt}
\begin{tabularx}{\textwidth}{>{\raggedright\arraybackslash}p{3.4cm}XXX}
\toprule
Свойство & zk-SNARK (Groth16) & zk-STARK & Bulletproofs \\
\midrule
Размер доказательства & сотни байт & десятки--сотни КБ & единицы КБ \\
Время проверки & константное, миллисекунды & логарифмическое & линейное \\
Доверенная установка & требуется; в PLONK универсальная & не требуется & не требуется \\
Постквантовая стойкость & нет (эллиптические кривые) & предполагается (хеш-функции) & нет \\
Применение & приватные платежи, rollup-системы & rollup-системы, доказуемые вычисления & доказательства диапазона \\
\bottomrule
\end{tabularx}
\end{table}

\section{Практическое применение}

\subsection{Приватные платежи и масштабирование блокчейнов}

Первое массовое применение zk-SNARK~--- криптовалюта Zcash (2016), где доказательство
подтверждает корректность транзакции при скрытых отправителе, получателе и сумме.
В версии Sapling (2018) время построения доказательства было снижено до единиц секунд,
что сделало приватные платежи практичными, а обновление на конструкцию Halo~2 (2022)
позволило отказаться от доверенной установки. В системе Monero с 2018~года применяются
Bulletproofs, что существенно сократило размер конфиденциальных транзакций.

Второе направление~--- zk-rollups, механизмы масштабирования, при которых тысячи
транзакций выполняются вне основной сети, а в неё публикуется лишь итоговое состояние с
доказательством корректности всех переходов. Основная сеть проверяет одно компактное
доказательство вместо повторного исполнения всех транзакций. На этом принципе построены
zkSync, Starknet, Scroll и Polygon zkEVM. Показательно, что нулевое разглашение здесь
используется не ради приватности, а ради \emph{краткости} доказательства.

\subsection{Анонимная аутентификация и цифровые удостоверения}

Естественное развитие классических протоколов идентификации~--- предъявление свойств
личности без раскрытия самой личности. Пользователь доказывает, что обладает выданным
государством удостоверением и что дата рождения в нём соответствует возрасту старше
18~лет, не сообщая ни даты рождения, ни номера документа, ни других атрибутов
(селективное раскрытие). Такой подход рассматривается в рамках европейской инициативы
цифрового кошелька личности и в системах проверки возраста, где сбор избыточных
персональных данных нежелателен по требованиям законодательства о защите данных.

\subsection{Защита авторских прав и подтверждение происхождения контента}

Для дисциплины, посвящённой защите авторских прав, наиболее интересны три сценария.

Во-первых, \emph{доказательство владения произведением без его раскрытия}: автор
публикует обязательство (например, хеш-обязательство на произведение и на встроенный
водяной знак) и при споре доказывает, что владеет прообразом, не публикуя само
произведение и не раскрывая параметры водяного знака. Это снимает известную проблему
стеганографических систем защиты авторских прав: предъявление ключа водяного знака в
суде фактически обесценивает сам знак, поскольку позволяет его удалить.

Во-вторых, \emph{подтверждение происхождения контента}. Развиваемые сегодня схемы
маркировки изображений и видео (в том числе сгенерированных нейросетями) опираются на
подписанные метаданные, однако любое преобразование файла~--- кадрирование, изменение
размера~--- подпись разрушает. Доказательство с нулевым знанием позволяет доказать, что
опубликованное изображение получено из подлинного оригинала применением разрешённого
набора преобразований, не раскрывая оригинала.

В-третьих, \emph{лицензионный контроль}: устройство или программа доказывает наличие
действующей лицензии, не передавая на сервер идентификатор пользователя или устройства,
что ограничивает возможности слежки со стороны правообладателя.

\subsection{Смарт-карты и встроенные системы}

Классические протоколы не потеряли актуальности там, где ресурсы ограничены. Механизмы
семейства Гиллу~--- Кискатра стандартизованы (ISO/IEC~9798-5) и применяются в
системах аутентификации устройств. Их привлекательность в том, что секрет никогда не
покидает карту и не передаётся терминалу даже в зашифрованном виде: компрометация
терминала не даёт противнику ничего, что можно было бы использовать повторно, в отличие
от схем с передачей пароля или PIN-кода.

\section{Ограничения и риски}

\begin{enumerate}
\item \textbf{Стоимость построения доказательства.} Асимметрия ресурсов принципиальна:
проверка занимает миллисекунды, а построение доказательства для сложного вычисления
требует секунд или минут и значительной памяти. Это ограничивает применение на
маломощных устройствах.

\item \textbf{Доверенная установка.} Для схем, требующих её, безопасность всей системы
зависит от корректности разовой церемонии, проверить которую задним числом невозможно.
Отсюда переход к прозрачным конструкциям.

\item \textbf{Ошибки в описании ограничений.} Доказывается ровно то, что записано в
системе ограничений. Недоопределённая программа порождает доказательство, которое
успешно проверяется, но не гарантирует того, что предполагал разработчик. Такие ошибки
обнаруживаются тяжело, а цена их высока.

\item \textbf{Ошибки реализации протокола.} Описанное выше слабое преобразование
Фиата~--- Шамира~--- характерный пример: библиотека проходит все функциональные тесты и
при этом допускает подделку доказательств.

\item \textbf{Побочные каналы.} Нулевое разглашение~--- свойство математического
протокола, а не его реализации. Время вычислений и энергопотребление устройства могут
раскрывать секрет, поэтому реализация на смарт-картах требует защиты от атак по
побочным каналам.

\item \textbf{Постквантовая перспектива.} Конструкции на эллиптических кривых теряют
стойкость при появлении масштабируемого квантового компьютера; схемы, опирающиеся
только на хеш-функции, считаются устойчивыми, но проигрывают по размеру доказательства.
\end{enumerate}

\sectioncenter{Заключение}

Доказательства с нулевым знанием прошли путь от теоретической конструкции, выглядевшей
парадоксально, до инженерной технологии массового применения. Все рассмотренные
классические протоколы~--- Фиата~--- Шамира, Гиллу~--- Кискатра, Шнорра~--- построены по
единой трёхшаговой схеме <<обязательство~--- вызов~--- ответ>> и различаются лишь
математической задачей, лежащей в основе, и способом уменьшения числа раундов. Их
развитие определялось практическими ограничениями смарт-карт, а преобразование
Фиата~--- Шамира связало эти протоколы с цифровой подписью, превратив интерактивный
протокол в подпись Шнорра и далее в применяемую сегодня схему EdDSA.

Современные конструкции сместили акцент с идентификации личности на доказательство
корректности произвольных вычислений. Практическим результатом стали приватные
платёжные системы, механизмы масштабирования блокчейнов, где нулевое разглашение
используется прежде всего ради краткости доказательства, и системы селективного
раскрытия персональных данных. Для защиты авторских прав технология даёт решение
давней проблемы: подтвердить владение произведением или параметрами водяного знака,
не раскрывая их и тем самым не обесценивая саму защиту.

Основные ограничения лежат не в математике, а в инженерии: стоимость построения
доказательства, риски доверенной установки, ошибки в описании доказываемых утверждений
и в реализации протоколов. Развитие идёт в сторону прозрачных и постквантово стойких
конструкций, что позволяет ожидать дальнейшего расширения области применения.

\sectioncenter{Список использованных источников}

\begin{enumerate}
\item Goldwasser, S. The Knowledge Complexity of Interactive Proof-Systems /
S.~Goldwasser, S.~Micali, C.~Rackoff // Proceedings of the 17th ACM Symposium on Theory
of Computing.~--- 1985.~--- P.~291--304.
\item Fiat, A. How to Prove Yourself: Practical Solutions to Identification and Signature
Problems / A.~Fiat, A.~Shamir // Advances in Cryptology~--- CRYPTO'86.~--- 1987.~---
P.~186--194.
\item Feige, U. Zero-Knowledge Proofs of Identity / U.~Feige, A.~Fiat, A.~Shamir //
Journal of Cryptology.~--- 1988.~--- Vol.~1, №~2.~--- P.~77--94.
\item Guillou, L. A Practical Zero-Knowledge Protocol Fitted to Security Microprocessor
Minimizing Both Transmission and Memory / L.~Guillou, J.-J.~Quisquater // Advances in
Cryptology~--- EUROCRYPT'88.~--- 1988.~--- P.~123--128.
\item Schnorr, C. Efficient Signature Generation by Smart Cards / C.~Schnorr //
Journal of Cryptology.~--- 1991.~--- Vol.~4, №~3.~--- P.~161--174.
\item Groth, J. On the Size of Pairing-Based Non-interactive Arguments / J.~Groth //
Advances in Cryptology~--- EUROCRYPT 2016.~--- 2016.~--- P.~305--326.
\item Ben-Sasson, E. Scalable, Transparent, and Post-Quantum Secure Computational
Integrity / E.~Ben-Sasson [et al.] // Cryptology ePrint Archive.~--- 2018.~---
Report 2018/046.
\item Bünz, B. Bulletproofs: Short Proofs for Confidential Transactions and More /
B.~Bünz [et al.] // IEEE Symposium on Security and Privacy.~--- 2018.~--- P.~315--334.
\item Ярмолик, С.~В. Криптография и охрана коммерческой информации: метод. пособие по
выполнению лаб. работ / С.~В.~Ярмолик, В.~Н.~Ярмолик.~--- Минск: БГУИР, 2010.~--- 32~с.
\item Шнайер, Б. Прикладная криптография. Протоколы, алгоритмы, исходные тексты на языке
Си / Б.~Шнайер.~--- М.: Триумф, 2002.~--- 816~с.
\end{enumerate}

\end{document}
